\chapter*{この問題集について}

この問題集に収録した問題および解説は全て、日本語の競技早押しクイズの問題を生成するAgent Skill \texttt{generate-quiz} (\url{https://github.com/Watson496/generate-quiz}) を用い、AIによって生成した。

\begin{description}[style=multiline,leftmargin=7\zw,labelwidth=6\zw]
  \item[作成日] \QuizCreationDate
  \item[使用ツール] \QuizGenerationTool
  \item[使用モデル] \QuizGenerationModel
\end{description}

\section*{generate-quizが生成する問題について}

\subsection*{題材、難易度、長さ}

generate-quizは題材を、\textbf{ファセット}と呼ぶ四つの選択軸で分類している。subjectは題材の分野、placeは題材にとって意味のある地域、timeは時代、typeは人物・作品・制度・出来事・概念など、解答対象がどのような種類のものかを表す。subject・place・timeの分類には国際十進分類法（UDC）の主題分類と共通補助表を、typeにはオントロジーに基づく分類を援用している。これらは資料を整理するための分類をそのまま示すものではなく、クイズの題材空間を区切る選択軸として用いている。subjectを中心とし、placeとtimeはその題材へ必要な場合に加わる補助的な限定である。

題材の選定では、同じ親分類に属する候補へ、問題群における自然な出現比率を表すweightを与える。基礎weightは、日本語圏における文化的共有度、多様なコミュニケーション場面での重要性、文化的背景知識としての機能から推定して与える。そのうえで、直近に選ばれたものと同じ分野、地域、時代、typeほどweightを下げ、問題が離れるにつれて基礎のweightへ戻る補正をかけてから重み付き乱択を行う。このため、特定のsubject・place・time・typeが短い間隔で続きにくい一方、最大weightの題材が常に選ばれるわけでもない。また、同じ命題の再利用を避け、固有性の高い人物・作品・出来事・組織等については、分類が異なっても数百問から数千問の尺度で再出現を抑えている。

解説の「題材選択」欄では、各ファセットを広い分類から狭い分類へ進む階層として表示する。placeまたはtimeが「限定なし」となっている場合、その問題に地域や時代の制約がないことを表し、資料不足や分類漏れを意味しない。四つの分類は問題を題材空間の中で位置づけるためのものであり、問題の優劣を表すものではない。

generate-quizが既定で生成する問題の難易度は、当該分野を継続的に学び始めて1～2年以内の人には基礎・概説知識として正答を期待できる一方、その分野を特に学んでいない一般的な日本語話者には通常は正答を期待しにくい水準としている。この「1～2年」は、大学1～2年生程度の専門知識という発想を、学問以外の分野にも適用できるように広げた目安である。実測正答率を示すものではなく、解答対象の知名度と、完成した問題文から正答するために必要な知識の双方に対する作問上の評価である。

問題文は80字付近を最も採用しやすくし、短い問題と長い問題も含みつつ、長い側にやや裾を引く分布になるよう調整している。各問を一定の字数範囲へ収める規則ではない。字数を合わせるためだけに周辺的な事実を足したり、成立に必要な情報を削ったりはしない。

\subsection*{問題文の構文}

generate-quizは問題文を、\textbf{SC型}または\textbf{OV型}の2つの構文に限定している。SC型は、疑問詞を解答で置き換えたときに「Nはxである」という平叙文になる構文であり、典型的には「～は何でしょう？」「～は誰でしょう？」の形を取る。対象そのものを尋ねるSC型では、同じ対象を正当に指す複数の名称が正答になり得る。一方、\textbf{OV型}は、解答を補うと「Nをxという」「Nをxと呼ぶ」という平叙文になる構文であり、典型的には「～を何というでしょう？」の形を取る。OV型では対象そのものではなく呼称を尋ねるため、必要に応じて「英語で」「正式名称で」「総称して」など、要求する名称の範囲を問題文で限定する。

名称を尋ねる問題が常にOV型になるわけではない。「正式名称は何でしょう？」のように、名称を対象としてSC型で尋ねる方が自然な場合もある。generate-quizは、対象そのものを同定させるのか、特定の名称・呼称を答えさせるのかと、採用する正答範囲から構文を選ぶ。SC型・OV型で自然に表せない題材を別の構文へ逃がすことはせず、答えさせる内容または題材を変更する。列挙問題や、年・数値・数量そのものを答える問題は、指定がない限り生成しない。

情報を担う節は、原則として、疑問詞へ続く核名詞句を修飾する部分に収める。文全体へ係る独立した前置き節、複数の文を並べる重文、前半を受けて問いへ転じる「ですが」構文は用いない。複数の独立した属性を入れる場合も、それぞれを核名詞へ係る連体修飾として構成する。連用中止は、並列・継起・理由・対立等の意味関係が実際に成立するときに限って用いる。このため、収録問題は一般的な競技クイズ問題集よりも構文の種類が意図的に狭く、各手掛かりが最終的に一つの解答対象または名称へ収束する形になっている。

この構文制約の下でも、読み上げられた順に意味を理解できることを求めている。独立した前フリとして働く連体修飾節は、後続名詞を選別する条件としてだけでなく、解答対象について事実を一つずつ加える叙述として自然に聞こえるようにする。後半の語によって前半の長い部分を大きく読み替えさせる構文や、修飾先・主体・因果関係を複数通りに読める構文は避けている。

手掛かりは、原則として知名度の低い情報を前に、高い情報を後ろに置く。ただし、日本語としての自然さを優先し、時系列、因果、係り受け、終端の落としを壊してまで機械的に並べ替えることはしない。

\subsection*{題材を代表する手掛かり}

generate-quizは問題文に、入門書・概説書・専門事典・博物館・学会等が、その対象を説明する際に中心に置く情報を用いる。著名な対象を、珍しい逸話、一時的なニュース、特殊な数値などで覆って人工的に難しくした問題は、原則として生成対象にしない。

中核的な情報であっても、それだけで適切な手掛かりになるとは限らない。generate-quizは、意味上ひとまとまりになる各手掛かりに、次の三つを個別に求めている。

\begin{itemize}
  \item \textbf{中核性}：その情報が、対象を理解するうえで中心的な位置を占めること。
  \item \textbf{準一意性}：現実的な比較範囲で、有力な近接候補を十分に絞れること。数学的に完全な一意性までは求めないが、同程度に有力な候補が残る情報は採用しない。
  \item \textbf{非重複性}：ほかの手掛かりの言い換えではなく、問題へ固有の情報を加えること。
\end{itemize}

これらは問題文全体に対する条件ではなく、各手掛かりに対する条件である。後半まで読めば答えが一意になる場合でも、単独では対象をほとんど絞れない前フリを残さない。前フリにも落としと同様に、中核性と準一意性を求める。落としは、解答対象と同じ抽象度・範囲で、対象が何であるかを端的に説明するものとしている。

また、ある事実が対象の一部分について成り立つことと、その事実で対象全体を特徴づけられることを区別している。例えば、著作中の著名な一節について正しい説明であっても、それだけで著作全体の代表的内容とはみなさない。一回の出来事を反復的な活動として述べたり、一例の存在を作風・傾向へ広げたりせず、資料が支える事実の種類と問題文の述語を一致させている。

\subsection*{問題としての非自明性}

問題文が事実として正しく、解答を一つに絞れても、知識問題として成立するとは限らない。generate-quizは、手掛かりと解答対象全体の間に、対象固有の非自明な対応があることを求めている。

例えば、「児童を支援する基金」という説明から「児童基金」という名称部分を作り、設立者の名前とつなぐだけで正式名を完成できる場合、実質的に問われている知識は設立者だけである。基金全体に固有の別の中核的情報がなければ、その基金を解答対象として選ぶ理由がない。このように、説明を複合語化・翻訳・同義置換するだけで解答名を作れる問題や、別の人物・対象を当てて自明な一般名詞を付けるだけの問題は生成しない。名称の構成が問題にならない題材では、対象全体と代表的な手掛かりとの非自明な対応を確認している。

\section*{generate-quizが出力する解説について}

generate-quizは各問題に、問題文、解答、手掛かりの選択、正誤判定、事実確認の根拠を読者が追跡するための解説を付ける。

\subsection*{解答と正誤判定}

「解答」には読み上げ用の代表解を一つ掲げ、必要に応じて読みや原語表記を添える。「別解」には、正式名、旧名、別題など、それ自体で同じ対象を完全に指定する別の名称を掲げる。姓のみ、名称の一部、表記・発音の軽微な揺れは、独立した別解ではなく「正誤判定基準」で扱う。

人名について「西洋人名は姓のみで○、東洋人名はフルネーム必須」といった機械的な規則は採用していない。その回答だけで、現実的に候補となるほかの人物・対象と区別し、想定した対象を十分に指定できるかを問題ごとに判断する。判定では、原則として現在の日本における標準的な表現を求める。一方、専門家コミュニティの間で正誤の合意が形成されていない表現は、明確に誤りとの合意もない限り正答として扱う。正答とも誤答とも確定できない回答には「もう一度」を求め、別対象を明示する誤った要素があれば×とする。問題文中の語と回答者の断片的な発話をつないで、回答者が発話していない名称を補うことはしない。

\subsection*{難易度と問題成立性}

「難易度評価」は、当該分野の初学者と分野外の一般層を分け、解答対象と完成した問題のそれぞれについて、既定または指定された難易度へ適合すると判断した理由を示す。これは事実関係の裏取りとは異なり、利用できる入門資料等を参考にした推定である。

「問題成立性の確認」は、実際の解答と代表的な手掛かりを挙げ、両者のどの対応を非自明な知識として問うのかを説明する。名称を容易に組み立てられないか、別の対象を当てるだけの問題へ縮約されないかも、その問題に即して扱う。抽象的な品質基準の合格宣言ではなく、何を知っていれば正答できる問題なのかを確認する欄である。

\subsection*{裏取り情報}

「裏取り情報」は、参考URLの一覧ではない。問題文が述べる内容を、異なる述語・関係ごとの命題に分け、次の三つを対応させている。

\begin{enumerate}
  \item その命題を生じさせる問題文の原文箇所
  \item 解答を補ったとき、その箇所から実際に読み取れる命題
  \item 命題の条件を支える出典の逐語引用と所在
\end{enumerate}

この対応により、出典に関連する語が見つかったというだけでなく、問題文の主体、対象、場所、時点、条件、因果関係等を、資料がどこまで支えるかを確認できる。複数の条件を含む命題では、一部の条件だけを立証する引用を命題全体の根拠にはしない。資料に合わせて解釈で問題文を修復するのではなく、資料が支える範囲を越えていれば問題文そのものを改める。

出典は、実在、作成主体、作成時期、原典か転載か、専門性、編集体制と、記述の射程や限定条件の両面から評価する。出典の信頼性は、原典・公式記録、専門機関、査読論文・専門辞典、専門家による解説、一般的なWeb解説の順を目安とする。資料数には固定の下限を置かず、転載された同じ記述を独立した複数の証拠とはみなさない。支持する資料だけでなく、異なる定義、例外、反例、対抗候補、時代や言語圏による用法の差も探す。

「情報選択の確認」では、各手掛かりについて、比較した対象の範囲、実際の有力な対抗候補、中核性、準一意性、ほかの手掛かりにはない寄与を示す。部分についての事実が真であることと、それで対象全体を特徴づける妥当性が異なる場合は、両者を分けて説明する。「解答の一意性」では、分類から近い対象を探すだけでなく、手掛かりの性質から逆向きに探した候補も含めて、解答との違いを示す。

「検証状況」は、根拠資料の該当本文を閲覧した「直接確認」、閲覧可能な二次資料が未閲覧の原資料を具体的に引用・報告した「間接確認」、関連性は推測できても命題を支える本文を確認していない「参照候補」を区別する。参照候補は証拠として扱わない。資料から結論を得る方法についても、出典への直接記載、確認済みの前提からの演繹、解釈・総合を区別する。問題全体の状態は、成立に必要な要素のうち最も弱いものに合わせる。「留保」には、定義差、異説、時点依存、間接確認、未確認事項、使用前の再確認事項を記す。「参考文献」には、問題文、解答、正誤判定、情報選択の検討に用いた資料を、本文中の参照箇所と対応が分かる形で掲げる。

これらの解説は、AI agentが生成した問題の正しさを保証するものではない。採用した事実、推論、評価、判定の根拠と限界を明示し、問題を使用する読者が再検討できるようにするためのものである。
